\documentclass[UTF8,zihao=-4]{ctexart}
\usepackage[a4paper,margin=2.5cm]{geometry}
\usepackage{booktabs,longtable,array,tabularx,enumitem,hyperref,xcolor}
\hypersetup{colorlinks=true,linkcolor=blue,citecolor=blue,urlcolor=blue}
\setlist[itemize]{leftmargin=2em}
\setlist[enumerate]{leftmargin=2em}
\title{从虚拟物品到数字自我：乙游角色资产、情绪价值与持续消费\\{\large Cluster C 论文写作支架（非最终正文）}}
\author{作者待填写}
\date{\today}
\begin{document}
\maketitle
\begin{center}\fbox{\parbox{0.9\textwidth}{\textbf{作者边界：} 本文档是写作支架，不是最终论文正文。所有正式段落、结论措辞和投稿文本必须由作者自行撰写和确认。}}\end{center}
\begin{abstract}
\noindent 摘要槽：作者在完成研究设计、数据收集和分析后填写。当前版本只保留研究对象、核心路径、方法和贡献的填写提示。
\end{abstract}
\section{研究定位}
\begin{itemize}
\item 主线：Cluster C，数字自我、虚拟商品与游戏消费。
\item 核心路径：虚拟角色资产的数字自我延伸 $\rightarrow$ 虚拟物品情感价值 $\rightarrow$ 角色依恋/品牌依恋 $\rightarrow$ 持续消费。
\item 研究对象：近3或6个月玩过国内女性向恋爱手游的成年玩家。
\item 分析单位：首选玩家-最喜欢角色关系。
\item 报告边界：横截面问卷只能支持关联和机制一致性，不能写成因果效应。
\end{itemize}
\section{引言写作骨架}
\begin{enumerate}
\item P1 现象入口：描述乙游作为女性向数字娱乐/情感消费场景；作者补行业事实或平台数据。
\item P2 理论入口：体验消费、享乐价值和品牌体验说明消费包含幻想、感受和乐趣。
\item P3 Cluster C 入口：角色、卡牌、账号资产和剧情记忆连接到数字自我和虚拟物品。
\item P4 研究缺口：现有文献分散，缺少乙游角色资产如何转化为情绪价值和持续消费的整合。
\item P5 研究问题：填写 population、exposure、mediator、outcome。
\item P6 方法概览：只写实际会做的数据和方法。
\item P7 贡献槽：等数据支持后再收紧贡献强度。
\end{enumerate}
\section{文献综述结构}
\subsection{体验消费与情绪价值}
主张槽：乙游消费可从体验消费和情绪价值理解。引用槽：Holbrook \& Hirschman (1982); Babin et al. (1994); Brakus et al. (2009)。
\subsection{数字自我与虚拟物品}
主张槽：角色、卡牌、账号资产可作为数字自我延伸或 loved objects。引用槽：Belk (1988, 2013); Ahuvia (2005)。
\subsection{角色依恋与品牌/IP依恋}
主张槽：玩家可能与角色和品牌/IP形成不同层次的关系。引用槽：Fournier (1998); Park et al. (2010); Batra et al. (2012)。
\subsection{社群与二创实践}
主张槽：社群参与和二创实践可能共同生产角色亲密感。引用槽：Kozinets (2002); Zhou et al. (2024)。
\section{构念与变量表}
\begin{longtable}{p{2.5cm}p{2.2cm}p{4.8cm}p{3.2cm}p{2.8cm}}
\toprule 构念 & 模型角色 & 工作定义 & 文献锚点 & 风险/下一步 \\ \midrule
\endfirsthead \toprule 构念 & 模型角色 & 工作定义 & 文献锚点 & 风险/下一步 \\ \midrule \endhead
数字自我延伸 & exposure & 玩家将虚拟角色、卡牌、账号资产、剧情记忆或收藏物视为自我表达、身份组成或个人经历延伸的程度 & Belk 1988；Belk 2013；Ahuvia 2005 & 可能与角色依恋高度相关；预测试后检查判别效度 \\
虚拟物品情感价值 & mediator & 玩家从乙游虚拟物品和角色资产中获得的愉悦、安慰、陪伴、审美和纪念性收益 & Babin et al. 1994；Holbrook \& Hirschman 1982；Brakus et al. 2009 & 可能与享乐价值和品牌体验重叠；用访谈确认乙游语境措辞 \\
角色依恋 & mediator\_or\_outcome & 玩家对特定乙游角色形成的稳定情感纽带和分离不适 & Park et al. 2010；Batra et al. 2012；Labrecque 2014 & 可能与拟社会关系重叠；决定是否将拟社会关系作为维度或相邻构念 \\
品牌/IP依恋 & outcome\_or\_parallel\_mediator & 玩家对游戏品牌、世界观、IP或运营方形成的情感纽带 & Park et al. 2010；Fournier 1998；Batra et al. 2012 & 乙游中玩家可能爱角色但不爱运营方；题项需明确对象 \\
持续消费意向 & outcome & 玩家未来继续游玩、继续付费或参与活动的主观意向 & customer engagement and brand loyalty literature & 不能等同于实际付费或留存；尽量加入行为区间题 \\
自报消费行为 & optional\_outcome & 玩家在指定时间窗内的自报付费、抽卡、月卡或活动消费行为 & methods review recommendation & 敏感题可能有缺失和低报；用区间题并允许不愿透露 \\
社群参与 & moderator\_or\_control & 玩家参与乙游社群、二创、讨论、委托、同担互动的程度 & Zhou et al. 2024；Kozinets 2002 & 可能是前因也可能是结果；模型中先作为控制或调节 \\
商业化反感/操纵感 & boundary\_condition & 玩家对限时活动、抽卡机制、情感话术或运营策略的压力和反感程度 & AI companion manipulation literature；methods review issue table & 敏感且可能引发防御性回答；放在后段并保护匿名性 \\
\bottomrule\end{longtable}
\section{研究命题槽位}
\begin{enumerate}
\item H1 槽：数字自我延伸与虚拟物品情感价值之间的正向关联。
\item H2 槽：虚拟物品情感价值与角色依恋之间的正向关联。
\item H3 槽：角色依恋与持续消费意向之间的正向关联。
\item H4 槽：虚拟物品情感价值和角色依恋的链式中介。
\item 边界变量槽：商业化反感/操纵感可能削弱情绪价值到持续消费意向的路径。
\end{enumerate}
\section{最小问卷题项}
\begin{longtable}{p{1.4cm}p{2.5cm}p{7.5cm}p{3.2cm}}
\toprule 题项ID & 构念 & 题项文本 & 量表/备注 \\ \midrule
\endfirsthead \toprule 题项ID & 构念 & 题项文本 & 量表/备注 \\ \midrule \endhead
DS1 & 数字自我延伸 & 我收藏或拥有的这个角色相关内容能体现我的个人偏好和身份 & 1-7 Likert 同意度；需根据具体乙游措辞调整 \\
DS2 & 数字自我延伸 & 这个角色或相关卡牌/剧情已经成为我游戏经历中很重要的一部分 & 1-7 Likert 同意度；预测试检查是否与角色依恋重叠 \\
DS3 & 数字自我延伸 & 看到这个角色相关内容时，我会想起自己投入过的时间、情绪或经历 & 1-7 Likert 同意度；记忆保存维度 \\
DS4 & 数字自我延伸 & 我会把这个角色相关的收藏视为“属于我”的重要数字资产 & 1-7 Likert 同意度；拥有感维度 \\
EV1 & 虚拟物品情感价值 & 收集或培养这个角色能给我带来愉悦感 & 1-7 Likert 同意度；享乐价值 \\
EV2 & 虚拟物品情感价值 & 这个角色相关内容能在情绪上安慰我或陪伴我 & 1-7 Likert 同意度；乙游语境核心 \\
EV3 & 虚拟物品情感价值 & 这个角色的视觉、声音、剧情或互动让我获得审美享受 & 1-7 Likert 同意度；品牌体验/审美体验 \\
EV4 & 虚拟物品情感价值 & 即使没有明显功能收益，我也觉得这个角色相关内容值得投入 & 1-7 Likert 同意度；区分功利价值 \\
CA1 & 角色依恋 & 我对这个角色有较强的情感连接 & 1-7 Likert 同意度；角色依恋核心 \\
CA2 & 角色依恋 & 如果之后很少再看到这个角色的新内容，我会感到失落 & 1-7 Likert 同意度；分离不适 \\
CA3 & 角色依恋 & 我会在意这个角色在剧情、活动或运营中的待遇 & 1-7 Likert 同意度；乙游语境 \\
CA4 & 角色依恋 & 相比游戏中的其他角色，这个角色对我更特别 & 1-7 Likert 同意度；角色特异性 \\
BA1 & 品牌/IP依恋 & 我对这款游戏或其IP有情感上的连接 & 1-7 Likert 同意度；对象为游戏/IP \\
BA2 & 品牌/IP依恋 & 即使暂时不玩，我仍会关注这款游戏或IP的新内容 & 1-7 Likert 同意度；持续关注 \\
BA3 & 品牌/IP依恋 & 这款游戏或IP对我来说不只是普通娱乐产品 & 1-7 Likert 同意度；品牌意义 \\
CI1 & 持续消费意向 & 未来一个月我愿意继续游玩这款游戏 & 1-7 Likert 同意度；继续游玩 \\
CI2 & 持续消费意向 & 未来如果有这个角色相关活动，我愿意继续参与 & 1-7 Likert 同意度；角色活动参与 \\
CI3 & 持续消费意向 & 未来如果经济条件允许，我愿意为这个角色相关内容付费 & 1-7 Likert 同意度；付费意向 \\
CI4 & 持续消费意向 & 我愿意向同好推荐这款游戏或这个角色相关内容 & 1-7 Likert 同意度；推荐意向 \\
PB1 & 自报消费行为 & 过去30天你在这款游戏中的付费金额区间是 & 0；1-30；31-100；101-300；301-1000；1000以上；不愿透露；敏感题 \\
PB2 & 自报消费行为 & 过去30天你是否为这个角色相关内容付费或消耗付费资源 & 是；否；不确定；角色相关行为 \\
PB3 & 自报消费行为 & 过去3个月你参与过多少次这个角色相关抽卡或活动 & 0；1-2；3-5；6次以上；不记得；回忆误差 \\
CP1 & 社群参与 & 我经常浏览这个角色或这款游戏相关的社群内容 & 1-7 Likert 频率；社群暴露 \\
CP2 & 社群参与 & 我会在社群中讨论、转发或创作这个角色相关内容 & 1-7 Likert 频率；主动参与 \\
CP3 & 社群参与 & 与其他玩家交流会增强我对这个角色的投入感 & 1-7 Likert 同意度；可能中介或调节 \\
MR1 & 商业化反感/操纵感 & 我有时觉得游戏利用我对角色的感情诱导消费 & 1-7 Likert 同意度；伦理边界 \\
MR2 & 商业化反感/操纵感 & 限时活动或抽卡机制有时会让我产生消费压力 & 1-7 Likert 同意度；付费压力 \\
MR3 & 商业化反感/操纵感 & 当角色情感内容和付费机制绑定过紧时，我会感到不适 & 1-7 Likert 同意度；情感商业化反感 \\
\bottomrule\end{longtable}
\section{分析计划槽位}
\begin{itemize}
\item 预测试：50-80人，检查题项理解、缺失率和初步判别效度。
\item 测量模型：报告 Cronbach alpha、CR、AVE、HTMT、载荷。
\item 主模型：SEM/PLS-SEM 检验数字自我延伸 -> 情绪价值 -> 角色依恋 -> 持续消费意向。
\item 替代模型：数字自我延伸 -> 角色依恋 -> 情绪价值 -> 持续消费意向。
\item 稳健性：区分付费/非付费玩家、高频/低频玩家；结果替换为自报消费行为。
\end{itemize}
\section{不支持或需谨慎的写法}
\begin{itemize}
\item 不要写：数字自我延伸导致持续消费。可写：数字自我延伸与持续消费意向呈关联。
\item 不要写：所有乙游玩家都会把角色视为自我延伸。可写：在样本或访谈材料中，部分玩家如此描述。
\item 不要写：平台商业化一定构成情感操纵。可写：情感商业化可能伴随商业化反感和付费压力。
\end{itemize}
\section{作者待决策事项}
\begin{enumerate}
\item 分析单位：玩家，还是玩家-角色关系？
\item 品牌/IP依恋是否保留，还是先聚焦角色依恋？
\item 是否收集真实或截图验证消费行为？
\item 商业化反感/操纵感放入主模型，还是作为讨论边界？
\end{enumerate}
\end{document}